% 本文件由 LaTeX 排版 Agent 根据有效 translation.json 自动生成。
% author=Irenaeus of Lyons; work_id=0134; title=从依勒内失传著作中辑录的残篇

\chapter{从\Person{依勒内}失传著作中辑录的残篇}

% structure: p0001 source=h1 target=omitted reason=page-title-already-rendered-by-parent

% structure: p0002 source=h2 target=section reason=relative-heading-rank; base=h2

\section{1}

凡誊写本书者，我因我们的主耶稣基督，并因祂荣耀的显现——当祂来审判活人死人时——恳请你，将你所誊写的与这本原稿对照，仔细校正；并请你同样将此恳词誊录于抄本之中。\par

% structure: p0004 source=h2 target=section reason=relative-heading-rank; base=h2

\section{2}

\Person{佛罗里诺}，这些见解（容我用温和的语气说）并非健全的教义；这些见解与教会不合，使其追随者陷于极度的不虔敬；这些见解，连教会以外的异端者也不曾敢于提出；这些见解，我们那些与宗徒交往的前任长老并没有传授给你。因为当我尚是孩童时，我在\Place{小亚细亚}见到你与\Person{波利卡普}在一起，你在宫廷中崭露头角，并竭力赢得他的赞赏。我对当时所发生之事的记忆比近期之事更为鲜明（因为童年的经历随灵魂一同成长，与之融为一体）；我甚至能描述那有福的\Person{波利卡普}通常坐着讲论之处——他的出入、他的一般生活方式与外貌，以及他向群众发表的讲论；也记得他如何讲述自己与\Person{若望}及其余见过主的宗徒的亲密交往，以及他如何追忆他们的话语。凡他从他们那里听来的关于主的事，无论是关于祂的奇迹还是教导，\Person{波利卡普}既从生命之言的目击者那里领受，便都按圣经和谐地讲述出来。这些事，借着天主赐予我的仁慈，我当时留心倾听，并珍藏在心里，而非纸上；而我如今不断地靠着天主的恩宠，在心中精确地反复思量。我能在天主面前作证，倘若那位有福的长老曾听见任何此类的事，他必会呼喊并掩耳，像他常做的那样说：\Quote{善良的天主啊，祢为我存留了怎样的时代，使我竟要忍受这些事？}他必会从听见那话的地方（无论坐着还是站着）立即逃离。这一点，从他写给邻近教会以坚定它们的书信，或写给某些弟兄以劝诫和鼓励的信中，也可以清楚看出。\par

% structure: p0006 source=h2 target=section reason=relative-heading-rank; base=h2

\section{3}

因为争论不仅仅关乎日子，也关乎禁食的形式本身。有些人认为自己应禁食一日，另一些人两日，还有些人更多，而另一些人则禁食四十日；他们将白昼和黑夜的时辰合计为他们的禁食日。这种守斋者之间的差异并非起源于我们的时代，而是远在我们的先辈时代，其中一些人或许因守斋不够精确，便因单纯或私人偏好而将习俗流传给后代。然而，尽管如此，所有这些人彼此和睦共处，我们也彼此保持和平。事实上，禁食的差异确立了共同信仰的和谐。而在你现今管理的教会中，先于\Person{索特尔}执政的那些长老——即\Person{阿尼塞特}、\Person{庇护}、\Person{希金}、\Person{泰勒斯福尔}和\Person{西斯克特}——他们自己并不如此守斋，也不允许与他们一起的人如此做。尽管如此，那些不这样守节的人，对那些从遵守此节的其他教区来的人，仍保持着和平的态度，尽管此种遵守与那些不遵行之人相比，更显得是明确的反对；从未有人因此事被逐出教会。相反，那些先于你的长老们，虽不遵守此俗，却将圣体送给其他教区中遵守此俗的人。当有福的\Person{波利卡普}在\Person{阿尼塞特}时代旅居\Place{罗马}时，尽管他们在某些其他点上发生了轻微争论，他们却立即彼此和好，不愿为此事在他们之间发生争执。因为\Person{阿尼塞特}既不能说服\Person{波利卡普}放弃他自己的守斋方式，因为这事一向由我们的主之门徒\Person{若望}和他所往来的其他宗徒如此遵守；而\Person{波利卡普}也不能说服\Person{阿尼塞特}按他的方式守斋，因为\Person{阿尼塞特}坚持他必须遵守前任长老的惯例。在此情况下，他们彼此共融；\Person{阿尼塞特}出于对\Person{波利卡普}的尊敬，在教会中让\Person{波利卡普}举行圣体礼；于是他们彼此和平分离，与整个教会——无论守此俗者还是不守者——保持和平。\par

% structure: p0008 source=h2 target=section reason=relative-heading-rank; base=h2

\section{4}

任何人只要有能力为邻人行善，却不去行，必将被视为与主的爱疏远的人。\par

% structure: p0010 source=h2 target=section reason=relative-heading-rank; base=h2

\section{5}

天主的旨意和能力是一切时间、地方、时代和一切本性的有效和预知的原因。旨意是理性灵魂的\OriginalTerm{理性}{\GreekText{λόγος}}，这理性在我们内，是灵魂所具有的、天赋自由行动的能力。旨意是渴求对象的理智，以及拥有智慧的、渴求所愿之物的欲望。\par

% structure: p0012 source=h2 target=section reason=relative-heading-rank; base=h2

\section{6}

既然天主是广大的、世界的建筑师、全能的，祂借着世界的建筑师和全能的意志创造了直达无限的事物，并以新的效果有力而有效地创造，以使所产生的一切事物的整个圆满得以形成，尽管它们先前并不存在——也就是说，凡不在我们观察之下的事物，以及我们眼前的事物。祂如此个别地容纳一切，并引导它们达到自己适当的结果，为了这个结果它们被召唤并产生，绝不改变成任何别的东西，除了它（结局）原本在本质上所是的。因为这是天主运作的特性：不仅进展到理智的无限，甚至超越我们心智、理性、言语、时间和地点以及每一个世代的能力；而且超越实体、圆满或完美。\par

% structure: p0014 source=h2 target=section reason=relative-heading-rank; base=h2

\section{7}

这种在主日不跪膝的习俗是复活的象征，借着它，我们因基督的恩宠从罪和死亡中得释放，死亡已在祂以下被消灭。如今这一习俗起源于宗徒时代，正如真福依勒内（\OriginalTerm{依勒内}{Irenaeus}），里昂（\OriginalTerm{里昂}{Lyons}）的殉道者兼主教，在他的论著\WorkTitle{论逾越节}中宣告的，在该书中他也提到了五旬节；在这一庆节我们也不跪膝，因为它与主日具有同等意义，原因已在上文述及。\par

% structure: p0016 source=h2 target=section reason=relative-heading-rank; base=h2

\section{8}

因为正如约柜内外部用纯金包裹，基督的身体也是如此纯洁和光辉；它内在由圣言装饰，外在由圣神庇护，以便从两者中清晰地显明诸本性的光辉。\par

% structure: p0018 source=h2 target=section reason=relative-heading-rank; base=h2

\section{9}

永远对配得的人说善言，从不说不配之人的恶语，我们也必将获得天主的荣耀和国度。\par

% structure: p0020 source=h2 target=section reason=relative-heading-rank; base=h2

\section{10}

确实，天主适合并符合其品性的是显示仁慈和怜悯，并为祂的受造物带来救恩，即使他们陷入毁灭的危险。经上说：\Quote{因为祂有赎罪（\OriginalTerm{赎罪}{propitiation}）。}\par

% structure: p0022 source=h2 target=section reason=relative-heading-rank; base=h2

\section{11}

基督徒的事业无非是时刻预备死亡（\OriginalTerm{预备死亡}{\GreekText{μελεπᾷν} \GreekText{ἀποθνήσκειν}}）。\par

% structure: p0024 source=h2 target=section reason=relative-heading-rank; base=h2

\section{12}

因此我们相信我们的身体也要复活。因为虽然它们归于腐朽，却不消亡；因为大地接受尸骸，保存它们，如同肥沃的种子混合在更肥沃的土壤中。同样，正如一粒赤裸的种子被播下，借着天主造物主的命令发芽，披上荣耀复活，但前提是它先死去、分解并与大地混合；由此看来，我们并非对身体的复活怀有虚妄的信仰。尽管它因原初的悖逆在指定的时间分解，但它被置于大地的熔炉中，好似要重新铸造；不再作为这腐朽的身体，而是纯洁、不再腐朽：这样每个身体都将恢复其自己的灵魂；当它披上这灵魂时，它不会经历忧伤，而会喜乐，永久保持纯洁状态，以正义的伴侣而非阴险的伴侣为伴，完全拥有属于它的一切，它将极其准确地领受这些；它不会领受与先前不同的身体，也不会摆脱痛苦或疾病，也不是作为荣耀的身体，而是如同它们离世时那样，在罪或义行中；它们如何，复活时也将如何披上；它们在不信中如何，也将按此被忠信地审判。\par

% structure: p0026 source=h2 target=section reason=relative-heading-rank; base=h2

\section{13}

当希腊人逮捕了基督徒慕道者的奴隶，然后对他们使用暴力，企图从他们那里得知基督信徒中某些秘密之事时，这些奴隶无话可说能满足折磨者的愿望，只听到主人说神圣的共融就是基督的身体和血，并以为那实际上是肉和血，就如此回答询问者。于是后者假设基督徒的实践确实如此，便将此事告知其他希腊人，并试图通过酷刑强迫殉道者圣桑克图斯（\OriginalTerm{圣桑克图斯}{Sanctus}）和圣布兰迪娜（\OriginalTerm{圣布兰迪娜}{Blandina}）承认这一指控。对这些人的回应，布兰迪娜以这些话语非常出色地回答道：\Quote{那些为了实践虔敬而不享用甚至允许他们吃的肉的人，怎能忍受这样的控告呢？}\par

% structure: p0028 source=h2 target=section reason=relative-heading-rank; base=h2

\section{14}

怎能说天主所创造的蛇原是哑而无理性的，却具有理性和言语呢？因为如果蛇本身有能力说话、辨别、理解并答复女人所说的话，那么就没有什么能阻止每一条蛇也这样做。但如果他们又说，这条蛇是按照天主的旨意和安排，用人的声音对厄娃说话，那么他们便使天主成为罪恶的制造者。邪恶的魔鬼也不可能将言语赋予一个无言的天性，从而从无中产生有；因为如果是那样，他就不会停止（为了迷惑人）藉着蛇、野兽和飞鸟与人交谈并欺骗他们。再者，作为一头野兽，它从哪里得知天主给那人独处的、隐秘的、甚至连女人自己也不知道的诫命呢？为什么它不攻击男人而攻击女人呢？如果你说它攻击她是因为两者中她较弱，[我回答说]，相反，她更强，因为她在违反诫命的事上似乎是男人的助手。因为她独自一人抵抗了蛇，并且在坚持了一会儿并抵抗之后，才被诡计所胜而吃了树上的果子；而亚当却毫不抵抗，也不拒绝，就吃了女人递给他的果子，这显示出极度的软弱和心智的懦弱。女人确实在战斗中败于魔鬼，是值得宽恕的；但亚当却不值得任何宽恕，因为他被一个女人战胜了——他本人亲自从天主那里领受了诫命。女人是从亚当那里听到诫命的，她轻视了它，或者因为她认为天主藉此说话是不相称的，或者因为她有怀疑，甚至可能认为这诫命是亚当擅自给她的。蛇发现她独自工作，所以能够与她单独交谈。看见她要么在吃树上的果子，要么不吃，他把那[禁]树上的果子摆在她面前。如果看见她正在吃，那么显然她是参与了可朽坏的身体。\BibleQuote{凡入口的，都排泄到厕所里去。}\BibleRef{玛 15:17}如果因此是可朽坏的，那么她显然也是会死的。但若是会死的，那么当然就没有诅咒；当天主的声音对那人说：\BibleQuote{你既是灰土，还要归于灰土，}\BibleRef{创 3:19}那时也没有[定罪]的判决，像事物的真实进程[现在并永远]进行那样。再者，如果蛇看见女人没有吃，那么它如何诱使从未吃过东西的她去吃呢？谁向这该死的杀人的蛇指出，天主对他们所说的死亡判决不会[立即]生效，当时祂说：\Quote{因为哪一天你吃了它，你必定要死？}不仅如此，而且随着[他们犯罪的]不受惩罚，那些直到那时还没有看见的人的眼睛将被打开？但伴随着所指的[眼睛]的打开，他们踏上了死亡之路。\par

% structure: p0030 source=h2 target=section reason=relative-heading-rank; base=h2

\section{15}

在古代，当巴郎以比喻说出这些话时，未被承认；如今，基督已出现并应验了它们，却仍不被相信。因此，[巴郎]预见此事并感到惊奇，喊道：\BibleQuote{哀哉！哀哉！天主实行这事时，谁能存活呢？}\BibleRef{户 24:23}\par

% structure: p0032 source=h2 target=section reason=relative-heading-rank; base=h2

\section{16}

再次为那跟随在旷野中被杀之人的一代解释法律时，他颁布了\WorkTitle{申命纪}；并非给他们一部不同于为他们祖先所规定的法律，而是综合了后者，以便他们通过听到发生在他们祖先身上的事，能全心敬畏天主。\par

% structure: p0034 source=h2 target=section reason=relative-heading-rank; base=h2

\section{17}

这些事预表了基督，祂被承认并被带入世界；因为祂在若瑟中被预表；然后按肉体，祂从肋未和犹大降生，作为君王和司铎；在圣殿中，祂被西默盎承认；通过则步隆，祂在外邦人中被相信，如先知所说：\BibleQuote{则步隆地；}\BibleRef{依 9:1}通过本雅明[即保禄]，祂因被传遍全世界而得荣耀。\par

% structure: p0036 source=h2 target=section reason=relative-heading-rank; base=h2

\section{18}

这并非没有意义；而是藉着十人的数目，他（基德红）可能显示耶稣作为他的助手，正如与他们所立的约所表明的。当他不愿与他们一起参与偶像崇拜时，他们将责任推到他身上：因为\Quote{耶鲁巴耳}意味着巴耳的审判座。\par

% structure: p0038 source=h2 target=section reason=relative-heading-rank; base=h2

\section{19}

\BibleQuote{你领\Person{若苏厄}——\OriginalTerm{若苏厄}{\GreekText{Ἰησοῦν}}——\Person{农}的儿子来。}\BibleRef{户 27:18}因为梅瑟带领人民出埃及是适宜的，但\Person{若苏厄}（耶稣）应带领他们进入产业。同样，梅瑟应如法律一样终止，但\Person{若苏厄}（\OriginalTerm{若苏厄}{\GreekText{Ἰησοῦν}}），作为圣言，以及那成为肉身的圣言（\OriginalTerm{成为肉身的圣言}{\GreekText{ἐνυποστάτου}}）的真实预表，应成为人民的传道者。此外，[适宜的是]梅瑟应以玛纳作为食物给祖先，而\Person{若苏厄}则给麦子；作为生命之初果，是基督身体的预表，正如圣经也宣告，当人民吃了从地上来的麦子时，主的玛纳就停止了。\BibleRef{苏 5:12}\par

% structure: p0040 source=h2 target=section reason=relative-heading-rank; base=h2

\section{20}

\BibleQuote{他就按手在他头上。}\BibleRef{户 27:23}\Person{若苏厄}的面容也因梅瑟按手而受光荣，但不及[梅瑟]的程度。既然他已获得一定程度的恩宠，[主]说：\BibleQuote{你将你的光荣分给他一些。}\BibleRef{户 27:20}因为[在此情形中]所赐之物并未停止属于赐予者。\par

% structure: p0042 source=h2 target=section reason=relative-heading-rank; base=h2

\section{21}

但他并不像基督那样藉着嘘气赐予，因为他不是圣神的泉源。\par

% structure: p0044 source=h2 target=section reason=relative-heading-rank; base=h2

\section{22}

\Quote{你不可与他们同去，也不可诅咒这人民。} \BibleRef{户 22:12} 他并没有暗示任何关于人民的事，因为他们都展现在他眼前，而是[指]预先指示的基督的奥秘。因为祂要按肉身从祖先出生，圣神预先给那人（巴郎）指示，免得他出于无知出去，对这人民发出诅咒。当然不是[他的诅咒]能产生任何违反天主旨意的效果；而是作为天主眷顾的展示，祂因他们祖先的缘故而施予他们。\par

% structure: p0046 source=h2 target=section reason=relative-heading-rank; base=h2

\section{23}

\Quote{他骑上了他的驴。} \BibleRef{户 22:22} 驴是基督身体的预像，所有的人都从劳苦中安息，如乘战车被承载其上。因为救主已背负了我们罪的重担。向巴郎显现的天使就是圣言本身，祂手中拿着一把剑，以表明祂从上而来的权能。\par

% structure: p0048 source=h2 target=section reason=relative-heading-rank; base=h2

\section{24}

\Quote{天主不像人。} \BibleRef{户 23:19} 祂由此表明，所有的人都确实有谎言之罪，因为他们从一事改变到另一事（\OriginalTerm{改变}{\GreekText{μεταφερόμενοι}}）；但天主并非如此，因为祂永远持续真实，成就祂所愿的一切。\par

% structure: p0050 source=h2 target=section reason=relative-heading-rank; base=h2

\section{25}

\Quote{要由主向米德杨复仇。} \BibleRef{户 31:3} 因为此人（巴郎）当他不再讲天主的神，而是违背天主的法律，设立关于淫乱的不同法律时，\BibleRef{户 31:16}当然那时不能算作先知，而是占卜者。因为那不遵守天主诫命的人，得到自己邪恶计谋的公正报应。\BibleRef{户 31:8}\par

% structure: p0052 source=h2 target=section reason=relative-heading-rank; base=h2

\section{26}

要知道每个人要么是空的，要么是满的。因为若没有圣神，他就不知道造物主；他没有接受耶稣基督——生命；他不知道在天上的父；若他不按照理性的指示、按照天上的法律生活，他就不是个明智的人，行为也不正直：这样的人是空的。另一方面，若他接受天主，天主说：\Quote{我要在他们中间居住，在他们中间行走，我要作他们的天主。} \BibleRef{肋 26:12}这样的人就不是空的，而是满的。\par

% structure: p0054 source=h2 target=section reason=relative-heading-rank; base=h2

\section{27}

因此，那引导三松手的小男孩，\BibleRef{民 16:26}预像了洗者若翰，他向人民显示对基督的信德。他们聚集的那房屋象征世界，其中居住着各种异教和不信的民族，向他们的偶像献祭。此外，两根柱子是两个盟约。于是，三松自己靠在柱子上，[表明]这点：当人民受教导时，认出了基督的奥秘。\par

% structure: p0056 source=h2 target=section reason=relative-heading-rank; base=h2

\section{28}

\Quote{天主的人说：\InnerQuote{掉在哪里？}他指给他那地方。他砍了一棵树，丢在那里，铁就浮了起来。} 这是一个记号，表明灵魂应通过木头被举升（\OriginalTerm{灵魂的举升}{\GreekText{ἀναγωγῆς} \GreekText{ψυχῶν}}），那木头是祂受苦的器具，祂能引导那些跟随祂升天的灵魂举升。这件事也指示了一个事实，即当基督的圣灵魂下降[到阴府]时，许多灵魂上升并在他们的身体中被看见。\BibleRef{玛 27:52} 就如木头是较轻的物体，沉入水中；而铁是较重的，却浮了起来：同样，当天主的圣言与血肉成为一体，通过一个身体与位格的结合，沉重和属地的[部分]被变为不朽的，在复活后，因天主性被举升到天上。\par

% structure: p0058 source=h2 target=section reason=relative-heading-rank; base=h2

\section{29}

\WorkTitle{玛窦福音}是写给犹太人的。因为他们特别强调基督[应为]达味后裔这一事实。玛窦自己也有更强烈的愿望[确定这一点]，特别费心向他们提供令人信服的证据，证明基督是达味的后裔；因此他以[叙述]祂的家谱开始。\par

% structure: p0060 source=h2 target=section reason=relative-heading-rank; base=h2

\section{30}

\Quote{斧子放在树根上，} \BibleRef{玛 3:10} 他说，催促我们认识真理，并藉着恐惧洁净我们，也预备[我们]在合适季节结出果实。\par

% structure: p0062 source=h2 target=section reason=relative-heading-rank; base=h2

\section{31}

请注意，借着比喻中的芥菜籽，天上的教义被表明，它像种子一样撒在世界里，如同在田里，[种子]具有固有的力量，火热而强大。因为全世界的审判者就这样被宣告，祂被隐藏在坟墓的地心三天，成为一棵大树，将祂的枝条伸展到地极。从祂发抽芽，十二宗徒，成为美好而多结果实的树枝，成为万民和天空飞鸟的庇护所；在这些枝条下，所有的人都如同鸟儿聚集到巢中一样，得以避难，并分享了从那枝条而来的有益健康的天上食物。\par

% structure: p0064 source=h2 target=section reason=relative-heading-rank; base=h2

\section{32}

若瑟夫斯说，当梅瑟在皇宫中被抚养长大时，他被选为将军攻打埃塞俄比亚人；获胜后，他娶了那国王的女儿，因为那女子出于对他的爱慕，将城交给了他。\par

既然这两人（\Person{亚郎}和\Person{米黎盎}）都对\Person{梅瑟}无礼，为什么只有后者受到了惩罚？\BibleRef{户 12:1} 首先，因为女人更应受责备，因为自然和法律都将女人置于男人从属的地位。或者，也许\Person{亚郎}在某种程度上可原谅，因为考虑到他是长兄，并享有大司祭的尊荣。此外，由于法律认为麻风病人不洁，而同时司祭职的起源和基础在于\Person{亚郎}，[主]并没有给他类似的惩罚，以免这污名沾染整个[司祭]种族；而是通过他姊妹的[例子]唤醒他的恐惧，并教导他同样的功课。因为\Person{米黎盎}的惩罚如此影响了他，以至于她一旦遭受惩罚，他就请求受伤害的[\Person{梅瑟}]，让他通过转求解除这痛苦。他也没有疏忽，而是立刻献上祈祷。于是，爱人类的主让他明白，祂不是作为审判者惩罚了她，而是作为父亲；因为祂说：\BibleQuote{如果她的父亲吐唾沫在她的脸上，她岂不应当羞愧吗？让她在营外关闭七天，然后才可以再进来。}\BibleRef{户 12:14}\par

% structure: p0067 source=h2 target=section reason=relative-heading-rank; base=h2

\section{33}

既然某些人，不知出于什么动机，从天主那里夺去了一半的创造能力，声称祂仅仅是物质中性质的原因，并主张物质本身是非受造的，现在让我们提出问题：凡在任何时候……是不变的。那么，物质是不变的。但如果物质是不变的，且不变之物在性质方面不发生任何变化，它就不构成世界的实体。因此，在他们看来似乎是多余的，即天主将性质附加于物质，因为物质不可能有任何改变，它本身是某种非受造之物。但进一步，如果物质是非受造的，那么它完全是按照某种性质而存在的，并且这性质是不变的，因此它不能接受更多的性质，也不能成为世界生成的材料。但如果世界不是由它而成的，那么[这种理论]就完全将天主排除在创造世界的能力之外。\par

% structure: p0069 source=h2 target=section reason=relative-heading-rank; base=h2

\section{34}

\BibleQuote{他浸了}，圣经说，\BibleQuote{七次在约旦河里。}\BibleRef{列下 5:14} 从前，\Person{纳阿曼}患麻风病，受洗后得洁净，这不是无故的，而是作为对我们的预表。因为正如我们在罪中成为麻风病人，我们藉着圣水和呼号主名，从我们过去的过犯中得到洁净；属灵上重生如同新生婴儿，正如主所宣告的：\BibleQuote{人若不藉着水和圣神重生，就不能进入天国。}\BibleRef{若 3:5}\par

% structure: p0071 source=h2 target=section reason=relative-heading-rank; base=h2

\section{35}

如果\Person{厄里叟}的尸体使死人复活，\BibleRef{列下 13:21} 那么当天主使人的死尸复活时，岂不更要领他们接受审判吗？\par

% structure: p0073 source=h2 target=section reason=relative-heading-rank; base=h2

\section{36}

那么，真知识在于对基督的理解，保禄称之为在奥秘中隐藏的天主智慧，就是十字架的道理，这是\BibleQuote{属血气的人不能领受的}\BibleRef{格前 2:14} 如果有人\BibleQuote{尝到了}\BibleRef{伯前 2:3} 他必不附和自高自大之人的争辩和诡辩，\BibleRef{弟前 6:4} 他们探究自己所不明白的事。\BibleRef{哥 2:18} 因为真理是纯朴的（\OriginalTerm{纯朴的}{\GreekText{ἀσχημάτιστος}}），并且\BibleQuote{这话离你甚近，就在你口里，在你心里}\BibleRef{罗 10:8} 正如同一宗徒所宣告的，对顺服之人易于领悟。因为它使我们像基督，如果我们经历了\BibleQuote{祂复活的能力和祂受苦的共融}\BibleRef{斐 3:10} 因为这就是教导的相似性和至圣的\BibleQuote{交给我们的信德}\BibleRef{犹 1:3} 是没有学问的人领受的，知识浅薄的人也教导了，不是\BibleQuote{专注重荒渺无凭的家谱}\BibleRef{弟前 1:4} 而是更致力于遵循正直的生活方式；免得因被剥夺了圣神而无法达到天国。因为真正的首要之事是舍己和跟随基督；那些这样做的人被引领达到成全，实现了他们老师的全部旨意，藉着属灵重生成为天主的子女，并成为天国的继承人；那首先寻求这些的人必不被撇弃。\par

% structure: p0075 source=h2 target=section reason=relative-heading-rank; base=h2

\section{37}

那些熟悉宗徒的次级规约（即在基督之下）的人，知道主在新约中设立了一种新的祭献，依照玛拉基亚先知的预言。因为，\BibleQuote{从日出之地到日落之处，我的名在外邦人中已受尊崇，在各处都向我的名献上馨香和纯洁的祭品；}\BibleRef{拉 1:11} 如若望在\BibleBook{若望}{默示录}中所宣告的：\Quote{香就是圣徒的祈祷。} 然后，保禄又劝勉我们\BibleQuote{将你们的身体献上，当作活祭，是圣洁的，是天主所喜悦的，这是你们合理的侍奉。}\BibleRef{罗 12:1} 并且又说：\BibleQuote{我们应藉着耶稣，常常以赞美的祭献，就是那以嘴唇颂扬他圣名之果，献于天主。}\BibleRef{希 13:15} 但这些祭献并非按照律法，律法的字据已被主取去，钉在十字架上；\BibleRef{哥 2:14} 而是按照圣神，因为我们必须在\BibleQuote{心神和真理}内朝拜天主。\BibleRef{若 4:24} 因此，圣体圣事的祭献不是肉体的，而是属神的；就此而言，它是纯洁的。因为我们向天主献上饼和祝福之杯，感谢祂命令大地结出这些果实来滋养我们。然后，当我们完成了祭献，我们就呼求圣神，使祂将这祭品显明出来：这饼成为基督的身体，这杯成为基督的血，好使领受这些预象的人获得罪的赦免和永生。因此，那些为纪念主而举行这些祭献的人，并不随从犹太人的观点，而是以属神的方式举行礼仪，他们将被称作智慧之子。\par

% structure: p0077 source=h2 target=section reason=relative-heading-rank; base=h2

\section{38}

宗徒们规定了，\BibleQuote{不可在食物、饮料或节期、月朔、安息日上论断任何人。}\BibleRef{哥 2:16} 这些争论从何而来？这些分裂从何而来？我们守节，却存着恶毒和邪恶的酵，割裂天主的教会；我们保存外表的事物，却丢弃了更美好的事物——信德和爱德。我们从先知的话中听到，这些节期和禁食不蒙主喜悦。\BibleRef{依 1:14}\par

% structure: p0079 source=h2 target=section reason=relative-heading-rank; base=h2

\section{39}

基督，在万世之前被称为天主子，在时期一满时显现，为要藉着祂的血洁净我们这些曾受罪恶权势辖制的人，将我们作为纯洁的儿子献给祂的父，只要我们顺从地接受圣神的惩戒。在时期的终结，祂将到来，铲除一切邪恶，并使万物和好，以便一切不洁得以终结。\par

% structure: p0081 source=h2 target=section reason=relative-heading-rank; base=h2

\section{40}

\BibleQuote{他找到一块驴腮骨。}\BibleRef{民 15:15} 应注意的是，在三松犯罪后，圣经不再提到他借助\BibleQuote{主的神临到他身上}这一公式所完成的事工。\BibleRef{民 14:6} 因为，按照圣宗徒的话，奸淫的罪是直接侵犯身体的，也涉及对天主圣殿的犯罪。\BibleRef{格前 3:16}\par

% structure: p0083 source=h2 target=section reason=relative-heading-rank; base=h2

\section{41}

这表明了仍处于不信中的外邦人所发动的对教会的迫害。但遭受这些事的他（三松）相信，对发动这场战争的人将有报复。但报复通过什么方式？首先，他投靠了不是感官所能认知的磐石\BibleRef{民 15:11}；其次，通过找到驴腮骨。而这驴腮骨的预像就是基督的身体。\par

% structure: p0085 source=h2 target=section reason=relative-heading-rank; base=h2

\section{42}

常对配得的人说好话，对不配的人从不说坏话，我们也必将达到天主的荣耀和国度。\par

% structure: p0087 source=h2 target=section reason=relative-heading-rank; base=h2

\section{43}

在这些事中，预言表明：子民因成为悖逆者，将被自己罪的锁链捆绑。但锁链自行断裂，则表明当他们悔改时，将再次从罪的桎梏中得到释放。\par

% structure: p0089 source=h2 target=section reason=relative-heading-rank; base=h2

\section{44}

一个灵魂在错误影响下，要接受相反的意见并非易事。\par

% structure: p0091 source=h2 target=section reason=relative-heading-rank; base=h2

\section{45}

\BibleQuote{他们用刀杀了彼珥的儿子巴郎。}\BibleRef{户 31:8} 因为他不再藉天主的神说话，反而设立了违背天主律法的另一条奸淫的律法，\BibleRef{默 2:14} 这人将不再被视为先知，而是占卜者。因为他没有持守天主的诫命，就得到了他邪恶计谋的公正报应。\par

% structure: p0093 source=h2 target=section reason=relative-heading-rank; base=h2

\section{46}

\BibleQuote{这世界的神；}\BibleRef{格后 4:4} 就是撒殚，对于不信者，他被称为神。\par

% structure: p0095 source=h2 target=section reason=relative-heading-rank; base=h2

\section{47}

若翰洗者的诞生终结了匝加利亚的哑疾。因为当声音从静默中发出时，他并没有加重父亲的负担；但正如不被相信时使他哑口无言，同样，当声音清晰地发出时，它释放了他的父亲——父亲曾被告知并生下了他。现在，这声音和燃烧的光\BibleRef{若 5:35}是圣言和光的先行者。\par

% structure: p0097 source=h2 target=section reason=relative-heading-rank; base=h2

\section{48}

因此，正如数字所表明的七十种语言，以及从分散中通过翻译合并为一种语言；同样，那方舟被宣称为基督身体的预像，这身体是纯洁无瑕的。因为正如那方舟内外都包裹纯金，同样，基督的身体也是纯洁而光辉的，内在由圣言装饰，外在由圣神庇护，以便两者同时彰显两种性质的辉煌。\par

% structure: p0099 source=h2 target=section reason=relative-heading-rank; base=h2

\section{49}

因此，藉着那早已产生之物，圣言现已赋予了一种解释。我们深信，在我们每人里面，可以说存在两个人：一个是隐藏的，另一个是显明的；一个是属肉体的，另一个是属神的；尽管二者的产生可以比作双胞胎。因为二者在世界面前只显为一个，因为灵魂在本性上并不先于身体；就其形成而言，身体也不先于灵魂：但二者是同时产生的；它们的滋养在于纯洁和甘甜。\par

% structure: p0101 source=h2 target=section reason=relative-heading-rank; base=h2

\section{50}

那时，对所有信而得生的人，将真正完成共同的大喜乐；在每人身上，复活的奥秘、不朽的希望、永恒国度的开始都将得以确认，那时天主已毁灭死亡和魔鬼。因为那从死者中复活的人性和肉身将不再死；但在转变为不朽、成为相似于神体之后，当天堂敞开时，[我们的主]充满荣耀地将它（肉身）呈献给圣父。\par

% structure: p0103 source=h2 target=section reason=relative-heading-rank; base=h2

\section{51}

然而，这些人的著作也许逃过了你们的注意，却落入了我们手中，我提请你们注意：为了你们的声誉，你们要把这些著作从你们中间驱逐出去，因为它们给你们带来了耻辱，因为这些著作的作者自夸是你们中的一员。因为对许多人来说，它们成了绊脚石，这些人单纯而毫无保留地接受这些著作，以为它们出自一位长老之手，却不知它们是对天主的亵渎。[试想]这些著作的作者，他通过它们不仅伤害了那些在[礼仪中]服事的人——这些人碰巧在思想上已准备好亵渎天主——而且也损害了我们中间的人，因为通过他的书，他给他们的心灵灌输了关于天主的错误教义。\par

% structure: p0105 source=h2 target=section reason=relative-heading-rank; base=h2

\section{52}

圣经论及基督承认：祂既是人子，同一位并非一个单纯的人；祂既是肉身，也是神体、天主圣言和天主。祂在末世由玛利亚诞生，同样也由天主出发，作为每受造物的首生者；祂饥饿时，却使[他人]饱足；祂口渴时，昔日却使犹太人得饮，因为\BibleQuote{盘石就是基督}\BibleRef{格前 10:4}自己：因此耶稣现今赐给信祂的人能力饮用涌出到永生的活水。\BibleRef{若 4:14}祂是达味之子，也是达味的主。祂源自亚巴郎，却先亚巴郎而存在。\BibleRef{若 8:58}祂是天主的仆人，也是天主子、宇宙的主。祂被人当众唾弃，却也将圣神嘘入祂的门徒。\BibleRef{若 20:22}祂忧愁时，却将喜乐赐给祂的子民。祂可被捉摸触摸，却以不可把捉的形态穿过那些试图伤害祂的人之中，\BibleRef{若 8:59}并毫无阻碍地穿过紧闭的门户。\BibleRef{若 20:26}祂睡觉时，却掌管海洋、风和风暴。祂受难时，却活着、赐予生命，并医治我们一切的软弱。祂死去时，却是死者的复活。祂在地上蒙羞，在天上却超越一切荣耀与赞美；祂\BibleQuote{虽然由于软弱被钉在十字架上，却由于天主的德能而活着}，\BibleRef{格后 13:4}祂\BibleQuote{降到了地下，}并\BibleQuote{升到了天上；}\BibleRef{弗 4:9}马槽已足容祂，祂却充满万有；祂曾死去，却永生永世。阿们。\par

% structure: p0107 source=h2 target=section reason=relative-heading-rank; base=h2

\section{53}

论及基督，法律、先知和福音作者都宣告：祂由童贞女所生，在木架上受难，并从死者中显现；祂也升了天，被圣父所荣耀，是永恒的君王；祂是完美的理智、天主圣言，在光明以前受生；祂是宇宙的创造者，与光一同，并是人的造主；祂是万有中的万有：诸圣祖中的圣祖，诸法律中的法律，诸司祭中的大司祭，诸君王中的统治者，诸先知中的先知，诸天使中的天使，诸人中的真人；在圣父内的圣子，在天主内的天主，直到永远的君王。因为正是祂与诺厄一同航行在[方舟]中，引导了亚巴郎；与依撒格一同被捆绑，与雅各伯一同流浪；是得救者的牧者，教会的新郎；革鲁宾的首领，天使军队的元首；出自天主的天主；出自圣父的圣子；耶稣基督；永远的君王。阿们。\par

% structure: p0109 source=h2 target=section reason=relative-heading-rank; base=h2

\section{54}

法律、先知和福音作者宣告：基督由童贞女所生，在十字架上受难，并由死者中复活，升天，被荣耀，永远为王。祂自身被称为完美的理智、天主圣言。祂是首生者，超越一切，人的创造主；万有中的万有：诸圣祖中的圣祖，诸法律中的法律，诸司祭中的司祭，诸君王中的元首，诸先知中的先知，诸天使中的天使，诸人中的真人；在圣父内的圣子，在天主内的天主，直到永远的君王。祂与若瑟一同被出卖，引导了亚巴郎；与依撒格一同被捆绑，与雅各伯一同流浪；与梅瑟一起，祂是领袖，关于百姓，是立法者。祂在先知中宣讲；由童贞女降生成人；生于白冷；被若翰接待，在约旦河受洗；在旷野受试探，并被证明是主。祂召集宗徒，宣讲天国；使瞎子复明，使死者复活；在圣殿中出现，却不为百姓所信；被司祭逮捕，带到黑落德面前，在比拉多面前被判罪；祂在肉身中显现，被悬于木架上，从死者中复活；显现给宗徒，之后升天，坐在圣父的右边，并被祂荣耀作为死者的复活。此外，祂是迷失者的救恩，黑暗中居民的光明，已生者的救赎；得救者的牧者，教会的新郎；革鲁宾的驭者，天使军队的领导者；出自天主的天主；耶稣基督我们的救主。\par

% structure: p0111 source=h2 target=section reason=relative-heading-rank; base=h2

\section{55}

\BibleQuote{那时，载伯德儿子的母亲，同她的儿子们前来，朝拜祂，请求祂一件事。}\BibleRef{玛 20:20}这些人当然不是没有理解力，那节经文所陈明的言论也不是没有意义：它们像序言一样预先陈述，与先前所阐释的那些要点有所一致。\par

\Quote{然后就近前来。}有时候，德行令人惊羡，不仅因其所彰显的本身，也因其彰显的时机。例如，葡萄或无花果的早熟果实，或任何果实，在其生长过程中，无人期待其成熟或完全发育；但即使某人察觉它尚有些许不完善，他也不会因此鄙视未熟的葡萄而认为无用，反而欣然采摘，因其出现得早；他并不考虑这葡萄是否有完美的甜味；反而立即因它早于其余而满足。同样，天主看见信徒拥有智慧，虽尚不完全，且只有微小的信德，祂也会原谅他们在这方面的不足，并不拒绝他们；反而慈祥地欢迎并接纳他们，如同初熟的果实，并尊重那印有德行的心思，即使尚未完全。祂容忍他们，视作收获前的先驱，并高度重视他们，因他们比其余人更早预备，仿佛预先占有了祝福。\par

因此，我们的先祖\Person{亚巴郎}、\Person{依撒格}和\Person{雅各伯}，应被尊崇于一切之上，因为他们确实为我们提供了如此早的德行榜样。有多少殉道者能与\Person{达尼尔}相比？请问，有多少殉道者能与\Place{巴比伦}的三位青年相媲美？虽然前者的记忆不如后者那般显著地呈现在我们面前。他们真正是初熟的果实，是后续丰硕收成的预兆。因此，天主命定其生平被记载，作为后来者的典范。\par

至于他们的德行如此被天主接纳，作为收成的初果，请听祂亲自宣告：\BibleQuote{正如葡萄，}祂说，\BibleQuote{我在旷野中找到了以色列，又如初熟的无花果，你们的祖先。}\BibleRef{欧 9:10}因此不要只因为亚巴郎相信而称他有福。你愿以惊叹的目光看亚巴郎吗？那就看当世界上六百人已被错谬污染时，他一人独宣虔诚。你愿达尼尔令你惊异吗？请看那\Place{巴比伦}城，在不敬神的嚣张与骄傲中昂首，其居民完全陷入各种罪恶。但达尼尔从深渊中浮出，吐出罪恶的盐水，喜跃投入虔诚的甘泉。同样，现在，关于\Person{载伯德}儿子们的母亲，不要仅惊叹她所说的话，还要惊叹她说这些话的时刻。因为她何时就近救主？不是在复活后，也不是在祂的名被宣讲后，更不是在祂的王国建立后；而是当主说：\BibleQuote{看，我们上耶路撒冷去，人子要被交给司祭长和经师；他们要杀死祂，第三天祂要复活。}\BibleRef{玛 20:18}\par

这些事是救主论及祂的受苦和十字架时所说的；祂向这些人预言了祂的受难。祂也没有隐瞒这受难将是最可耻的，出自司祭长之手。然而这妇人将祂受苦的救恩计划作了另一种解释。救主预言死亡，而她求得不朽的荣耀。主宣告祂必须站在不敬神的审判者面前受审；但她不顾那审判，却向审判者请求：\BibleQuote{准许，}她说，\BibleQuote{我这两个儿子，一个坐在祢右边，一个坐在祢左边，在祢的荣耀里。}一方面提及受难，另一方面却被理解为王国。救主谈论十字架，而她所关注的是不容受苦的荣耀。因此，正如我已说的，这妇人值得我们惊叹，不仅因她所求的，也因她提出请求的时机。\par

她确实受苦，不仅作为虔诚者，也作为妇人。因为，受祂话语的教导，她思考并相信此事必将实现：基督的王国要在荣耀中昌盛，以其辽阔遍行世界，并因虔诚的宣讲而扩展。她明白了（事实正是如此），那以卑微形像显现的祂，已交付并接受了所有应许。我将在另一场合，当我论及这谦卑时，探究主是否拒绝了关乎祂王国的请求。但她想，当天使显现、祂受天使服侍并接受整个天军的服务时，她将不再拥有同样的自信。因此，她将救主带到僻静处，恳切地向祂求那些超乎一切人性的事物。\par

\SourceRecord{Translated by Alexander Roberts.；Ante-Nicene Fathers；Vol. 1.；Edited by Alexander Roberts, James Donaldson, and A. Cleveland Coxe.；Buffalo, NY: Christian Literature Publishing Co.,；1885.；<http://www.newadvent.org/fathers/0134.htm>.}
